% Finite scale induction behind the endpoint price-duality estimate.
\section{The finite price recurrence}
\label{sec:ex-lattice-trace-price-recurrence}

We give the scale induction that turns the paid direction and the local
one-step response estimate into the endpoint inequality used in
Section~\ref{sec:ex-lattice-trace-coffer-supplier}. The one-step response
estimate is derived below from the finite-volume local insertion, local
uniqueness, and block-size bounds. The exact source declarations and their
remaining scale inputs are recorded in the accompanying evidence file.

Fix a reference degree \(d\), a selected scale history \(H\), and integers
\(m\leq k\). Write
\[
 \Gamma_j=\exp\!\left(\frac{S_k-S_j}{8}\right),\qquad
 \kappa_j=\sqrt{\operatorname{monitor}(j)},\qquad
 I_j=\frac1{\Gamma_j\kappa_j}\quad(m\leq j\leq k).
\]
Here \(S_j\) is the scale coordinate of the history. For \(m\leq j<k\),
let \(L_j\) be the block side length at step \(j\), let
\[
 P_j=C(J(j+1)),\qquad
 b_j=\frac1{10K(J(j+1))},\qquad
 \rho_j=\operatorname{rate}(L_j),\qquad
 s_j=\frac{b_j}{\kappa_{j+1}},
\]
where \(J(j)\) is the stage at rung \(j\), \(C\) is the selected
coefficient-cost profile, \(K\) is the fixed positive worth profile, and
\(\rho_j>0\) is the block remainder-cost factor. For each rung \(j\), write \(c_j\geq1\) for the selected column allowance of the history at that rung. It is a real parameter of the selected stage data, not the cardinality of a coefficient index set; \(c_k\) is its value at the terminal
rung, matching the notation of the Coffer estimate. Put \(p_j=L_j^{1/8}\), and define
\begin{equation}
 H_j=\Gamma_j^{-1}+\frac{P_jL_j^d+\rho_j}{\Gamma_{j+1}},\qquad
 \varepsilon_j=\frac{I_{j+1}}{10P_jH_j}.
 \label{eq:ex-lattice-price-local-scales}
\end{equation}
All quantities in the denominators are positive.

The scale and budget choices give, at every step,
\begin{equation}
 \kappa_{j+1}\leq2\kappa_j,\qquad
 1\leq P_j,\qquad 4P_j\leq p_j,\qquad
 2P_jb_j\leq\frac13,\qquad
 \Gamma_j=\Gamma_{j+1}p_j.
 \label{eq:ex-lattice-price-budget}
\end{equation}
Consequently
\begin{equation}
 P_jI_j\leq\frac12 I_{j+1}.
 \label{eq:ex-lattice-price-contraction}
\end{equation}
Indeed, the ratio of the left side to \(I_{j+1}\) is
\(P_j\kappa_{j+1}/(p_j\kappa_j)\), which is at most \(2P_j/p_j\leq1/2\).

For a coefficient \(x\) at rung \(j\), write \(x^{\circ}\) for its
projection to nonconstant coordinates, and let \(\mathcal L_j(x)\) be the
coefficient produced by the local block insertion. The diagonal weight map
\(\mathcal V_{L_j}\) multiplies a grade-\(r\) coordinate by
\(L_j^{d-r}\). Thus \(\|\mathcal V_{L_j}(x^{\circ})\|
\leq L_j^d\|x^{\circ}\|\).
The one-step local response estimate, for a paid comparison with adjacent coefficients
\(x_j,x_{j+1}\) and nonnegative remainder-cost bounds \(r_j,r_{j+1}\), is
\begin{equation}
 \begin{split}
 \|x_j-\mathcal L_j(x_{j+1})\|
 &\leq s_j\bigl(\|x_{j+1}^{\circ}\|+r_{j+1}\bigr)\\
 &\quad+\varepsilon_j\Bigl(r_j+P_j\|\mathcal V_{L_j}
                (x_{j+1}^{\circ})\|+\rho_jr_{j+1}\Bigr).
 \end{split}
 \label{eq:ex-lattice-price-one-step}
\end{equation}
Here \(\|x\|\) is the coefficient \(\ell^1\)-norm, the sum of the absolute
values of all coordinates including the unit coordinate, and \(x^\circ\)
deletes exactly the unit coordinate. In the selected construction, under
its standing scale hypotheses (in particular the stage at rung \(j\) is at
least \(d+3\)), the estimate
holds for all sufficiently large period indices \(n\), with a threshold
depending on \(j\) and on the constant \(\varepsilon_j\), and then uniformly
in the volume index \(N>j\) and in the pair of adjacent fields, provided
both fields belong to the effective representations of the same selected
history at the rungs \(j\) and \(j+1\) for these \(N\) and \(n\), have the
required order-zero format and remainder bounds at the normalizations
\(\sqrt{\operatorname{forcedAllowance}_j}\) and
\(\sqrt{\operatorname{forcedAllowance}_{j+1}}\) with costs
\(r_j,r_{j+1}\), satisfy the exact density transfer (the rung-\(j\) density
is \(L_j^d\) times the pushed rung-\((j+1)\) density), and have marked
sites related by the scale map. The separate
local insertion estimate is also
\begin{equation}
 \|x_{j+1}\|\leq P_j\|\mathcal L_j(x_{j+1})\|.
 \label{eq:ex-lattice-price-spends}
\end{equation}
We explain the finite-volume comparison producing the second line of
\eqref{eq:ex-lattice-price-one-step}. Let \(\mathcal Q_d\) be the finite
family of coefficient-recovery probes, with probe prices \(\pi(\tau)\) for \(\tau\in\mathcal Q_d\). Put
\[
 M_d=1+\sum_{q=0}^{d}2^q(q+1)^q .
\]
This is the polarization constant: if two functions are of class \(C^d\) near the identity and all ordered directional derivatives of orders \(0,\ldots,d\) of their difference at the identity, along normalized directions, are bounded by \(c\), then for every probe \(\tau\) the two probe values, in which the functions are composed with the logarithmic-link test curve of \(\tau\), differ by at most \(\pi(\tau)M_dc\). The factor \(M_d\) pays for expressing the mixed derivatives that occur in the chain rule through the available ordered derivatives, and \(\pi(\tau)\), defined below, pays for the composition with the test curve. Let \(T_d\) be the finite coefficient reconstruction constant for this probe family. Thus for any
coefficient difference \(z\),
\[
 \|z\|\leq T_d\sum_{\tau\in\mathcal Q_d}|R_\tau(z)|.
\]
A probe of order at most \(d\) annihilates a regular source term whose
support is nonwrapping and whose format has the required vanishing of its
first \(d\) local variations. Let \(B_0\) be the support radius and
\(\sigma>0\) the support exponent of the local representation, and define
\[
 M_{B_0}(X,W)=e^{\sigma\operatorname{cover}_{B_0}(X)}W.
\]
For each of the two coefficients, the regular part has a countable
presentation by regular terms with origin in their supports and with total
\(M_{B_0}\)-mass bounded by its remainder budget. Each term has a common
healthy chart. If its purchased sup-bound is \(u_a\), then on this chart
every ordered directional derivative of orders \(q=0,\ldots,d\) has norm at
most \(u_a\), after the local variation prices have been factored out.
The chart is large enough for each recovery probe and its local stencil to
fit, and the admissible local presentation has the required order-\(d\)
vanishing on each nonwrapping term. These are the finite recovery and
regularity hypotheses used below; a value bound alone would not imply the
probe estimate. Their signed union is still absolutely summable, even when
some supports overlap.

Let \(\mathcal I_0\) be the indices of nonwrapping terms, and take a finite
subset \(D\subset\mathcal I_0\). The probe applied to the presented
coefficient equals the corresponding derivative of its local presentation;
this is the coefficient-readout identity.
On the common quiet chart, the forced contributions vanish, and the
prescribed masks replace the local masks. Thus each density is the same
nonzero exponential action factor times its presented coefficient plus its
regular series. Equality of the two densities lets us cancel this common
factor and gives the signed difference of the presented coefficients as
the sum of all signed regular terms: writing \(g\) for the local
presentation of the coefficient difference and \(f_a\) for the signed
regular terms,
\begin{equation}
 g(U)=\sum_a f_a(U)
 \label{eq:ex-lattice-price-presented-sum}
\end{equation}
for every configuration \(U\) in the quiet set \(\mathcal D\), the set of
configurations all of whose plaquette variables are strictly below the
threshold. The set \(\mathcal D\) is open and contains the identity
configuration, every support is healthy on it, and the series converges
absolutely there because \(|f_a|\leq u_a\) and \(\sum_au_a<\infty\).

The probe estimate needs the ordered derivatives of \(g\), so the identity
\eqref{eq:ex-lattice-price-presented-sum} must be differentiated. We claim
that for every \(U\in\mathcal D\) and every ordered list
\((A_1,\ldots,A_q)\) of normalized directions with \(q\) at most the order
of the budget,
\begin{equation}
 \partial_{A_1}\cdots\partial_{A_q}g(U)
   =\sum_a\partial_{A_1}\cdots\partial_{A_q}f_a(U),
 \label{eq:ex-lattice-price-derivative-sum}
\end{equation}
where the series converges absolutely. The proof is by induction on \(q\);
the case \(q=0\) is \eqref{eq:ex-lattice-price-presented-sum}. For the
induction step, let \(W(t)\) be the increment curve through \(U\) in the
direction \(A_1\). Since \(\mathcal D\) is open, \(W(t)\in\mathcal D\) for
\(|t|<r\) with some \(r>0\). On this interval the functions
\(v_a(t)=\partial_{A_2}\cdots\partial_{A_q}f_a(W(t))\) are differentiable
by regularity of the terms on healthy configurations, their derivatives are
\(\partial_{A_1}\cdots\partial_{A_q}f_a(W(t))\), and these are bounded by
\(u_a\) uniformly in \(t\), with \(\sum_au_a<\infty\). The series
\(\sum_av_a(0)\) converges. The theorem on termwise differentiation of a
series whose derivatives are dominated by a summable sequence on an
interval therefore shows that \(\sum_av_a\) is differentiable at \(0\) with
derivative \(\sum_av_a'(0)\). By the induction hypothesis, applied at
every point \(W(t)\), \(|t|<r\), the sum \(\sum_av_a(t)\) equals
\(\partial_{A_2}\cdots\partial_{A_q}g(W(t))\); differentiating at \(t=0\)
gives \eqref{eq:ex-lattice-price-derivative-sum}.

Now fix the finite set \(D\) and let \(G_D=\sum_{a\in D}f_a\). The finite
sum \(G_D\) has zero probe: every term has the
vanishing property just stated, and a finite sum may be differentiated
termwise. By \eqref{eq:ex-lattice-price-derivative-sum} at the identity
configuration, every ordered derivative of \(g-G_D\) of order at most
\(d\) along normalized directions is the sum of the corresponding
derivatives of the omitted nonwrapping terms and of the wrapping terms, so
it is bounded by
\(\sum_{a\in\mathcal I_0\setminus D}u_a+\sum_{a\notin\mathcal I_0}u_a\).
The polarization estimate stated with the definition of \(M_d\) then bounds
the difference between the probe of the full presented
field and the zero probe of the truncation \(G_D\). This estimate is
uniform in finite \(D\), because the
local section has order-\(d\) regularity and the source terms have summable
sup-bound mass. If the probe has order \(r\), the chain rule expands its
local link variation into ordered partitions of \(\{1,\ldots,r\}\). Let
\(\lambda_{\tau,q}\) be the price of a local logarithmic variation of
order \(q\); the finite sum of the products
\(\prod_{A\in\mathcal P}\lambda_{\tau,|A|}\) over these partitions,
with the zeroth-order term included, is \(\pi(\tau)\). Supremum bounds
on the omitted terms and the triangle inequality therefore give
\[
 |R_\tau(P.x-Q.x)|\leq\pi(\tau)M_d\left(
   \sum_{a\in\mathcal I_0\setminus D}u_a
   +\sum_{a\notin\mathcal I_0}u_a\right),
\]
where \(u_a\) is the sup-bound of the corresponding signed regular
term. The remainder certificates give \(\sum_a M_{B_0}(X_a,u_a)\leq a+b\),
so \(\sum_a u_a\) is finite as well. For the wrapping terms, the elementary support inequality
\[
 10\bigl(h(4B_0+2)+2B_0\bigr)<n,
 \quad 0\in X,\quad X\text{ wrapping},\quad W\geq0
 \quad\Longrightarrow\quad
 W\leq e^{-\sigma h}M_{B_0}(X,W)
\]
follows from the lower bound on the cover radius of a wrapping support.
Absolute summability of \(\sum_au_a\) justifies the tail limit and permits
the finite sets \(D\) to increase
through all of \(\mathcal I_0\). The nonwrapping
sup-bound tail tends to zero, and the wrapping sup-bound sum is at most
\(e^{-\sigma h}(a+b)\); taking this limit gives the probe estimate
\[
 |R_\tau(P.x-Q.x)|\leq\pi(\tau)M_d e^{-\sigma h}(a+b).
\]
This is the truncation-and-tail argument. The probe is applied termwise
only to the finite truncations, where the vanishing property is available;
the infinite collection of omitted and wrapping terms enters only through
the derivative identity \eqref{eq:ex-lattice-price-derivative-sum} and the
resulting sup-bound.

The coefficient recovery inequality is
\[
 \|z\|\leq T_d\sum_{\tau\in\mathcal Q_d}|R_\tau(z)|.
\]
It follows from finite-dimensionality: the probes jointly separate
coefficients, so their maximum is a norm on the coefficient space; compactness
of its unit sphere gives a finite comparison constant, and a maximum is
bounded by the sum. Summing the probe estimate over the finite family gives
\begin{equation}
 \|P.x-Q.x\|\leq \mathfrak C_d e^{-\sigma h}(a+b),\qquad
 \mathfrak C_d=T_dM_d\sum_{\tau\in\mathcal Q_d}\pi(\tau).
 \label{eq:ex-lattice-local-uniqueness}
\end{equation}
Take the torus side length \(n\) so that
\(10(h(4B_0+2)+2B_0)<n\). Then every probe fits in the region of
nonwrapping supports. This proves the required local uniqueness estimate
from coefficient recovery, local vanishing, the exponential support tail,
and absolute summation.

For the insertion itself, first decompose the next-level density into its
primary term, its countable regular terms, and its countable forced terms.
The primary term has only finitely many coefficient coordinates.
The regular series is absolutely summable by its sup-bound mass budget. The
forced series is absolutely summable after division by the positive factor
\(y=\sqrt{\operatorname{forcedAllowance}_{j+1}}\), using the forced part
of the same remainder certificate. Write the regular mass sum as \(a\) and the forced mass sum as \(yb\),
where \(a,b\geq0\) and \(a+b\leq r_{j+1}\) by the remainder
certificate. For a regular term of volume mass \(m_a\), the paid atom
costs are \(s_jm_a\) and \(\rho_jm_a\); for a forced term of mass
\(m_a\), they are \(s_jm_a/y\) and \(\rho_jm_a/y\). Thus the
combined regular and forced costs are bounded by \(s_jr_{j+1}\) and
\(\rho_jr_{j+1}\), respectively. The primary coefficient atom contributes the finite
coordinate costs \(s_j\|x_{j+1}^{\circ}\|\) and
\(P_j\|\mathcal V_{L_j}(x_{j+1}^{\circ})\|\). Summing first over the
finite coefficient basis, then over the countable regular and forced lists,
is legitimate because their total costs and coefficient norms are summable.
The interchange of the countable sums with the integral needs a bound
for the factors multiplying the atoms in the density, and this is not a
consequence of compactness of the link variables alone. Two properties of
the fine effective representation are used; both are part of its
representation certificate and are inputs here. First, its effective
density is measurable and bounded. Second, its vacuum activities \(K_j\),
\(j\in\mathbb N\), have a summable budget: there are \(a_j\geq0\) with
\(\sum_ja_j<\infty\) and \(|K_j(U)|\leq a_j\) for all \(U\). The multiplier
that dresses the \(i\)-th forced insertion is the sum, over the finite
admissible index sets \(b\ni i\), of \(\prod_{j\in b\setminus\{i\}}K_j(U)\).
Erasing the index \(i\) is injective on such sets, so every dressed
multiplier is bounded, uniformly in \(i\) and \(U\), by
\[
 C_{\mathrm{gas}}=\sum_{b\subset\mathbb N\ \mathrm{finite}}
     \prod_{j\in b}a_j\leq\exp\Bigl(\sum_ja_j\Bigr)<\infty.
\]
Hence a summable majorant \(\beta_i\) for the forced insertions gives the
summable majorant \(C_{\mathrm{gas}}\beta_i\) for their dressed versions,
and the same absolute summability justifies rearranging the sum over
marked index sets into the sum of dressed insertions. The absolute
integrand of the countable atomic
sum is therefore bounded by the bound of the effective density times a
summable majorant. This is
an integrable bound on the normalized finite-volume Haar space, so Fubini
preserves the density identity.
Thus the countable atomic sum has the pushed density. Denote its actual
coefficient by \(\widetilde x_j\). The coordinate leads sum to
\(\mathcal L_j(x_{j+1})\), and the paid insertion estimate gives
\(\|\widetilde x_j-\mathcal L_j(x_{j+1})\|\leq
s_j(\|x_{j+1}^{\circ}\|+r_{j+1})\). Its coarse remainder budget is at
most \(P_j\|\mathcal V_{L_j}(x_{j+1}^{\circ})\|+\rho_jr_{j+1}\).
This is the content of the countable summation argument: dominatedness and
absolute summability justify both the density identity and the interchange
of summation with integration.

Compare this intermediate coefficient with the actual coefficient at rung
\(j\) using the local uniqueness estimate
\eqref{eq:ex-lattice-local-uniqueness}. Their local densities agree by the
paid comparison's density identity and the insertion's pushed-density
identity. Choose \(h\) so that
\(\mathfrak C_de^{-\sigma h}\leq\varepsilon_j\), then choose the torus
side long enough for the displayed support inequality. The triangle
inequality gives exactly \eqref{eq:ex-lattice-price-one-step}, uniformly
over the allowed fields.

For \eqref{eq:ex-lattice-price-spends}, coefficient coordinates have grades
\(r=0,1,\ldots,d\), so each weight \(L_j^{d-r}\geq1\) when \(L_j\geq1\).
Consequently the weighted coordinate norm dominates the original norm. The legal block cap
then controls that weighted norm by \(P_j\) times the norm after
insertion. This is the size estimate for the admissible leading map; it is
not merely the inequality \(L_j\geq1\).

\begin{lemma}[One-step propagation]
\label{lem:ex-lattice-price-one-step}
Suppose the adjacent data satisfy the local hypotheses above and, for some
\(t\geq0\),
\[
 r_j+\|x_j\|\leq\frac{t}{\Gamma_j},\qquad
 r_{j+1}+\|x_{j+1}\|\leq\frac{t}{\Gamma_{j+1}},\qquad
 \|x_j\|\leq tI_j.
\]
Then
\[
 \|x_{j+1}\|\leq tI_{j+1}.
\]
\end{lemma}

\begin{proof}
The projection to nonconstant coordinates is contractive, so the paid-cone
bound gives \(\|x_{j+1}^{\circ}\|+r_{j+1}\leq
\|x_{j+1}\|+r_{j+1}\leq t/\Gamma_{j+1}\). It also gives
\[
 r_j+P_j\|\mathcal V_{L_j}(x_{j+1}^{\circ})\|
       +\rho_jr_{j+1}
 \leq t\left(\Gamma_j^{-1}
        +\frac{P_jL_j^d+\rho_j}{\Gamma_{j+1}}\right)=tH_j.
\]
Apply \eqref{eq:ex-lattice-price-one-step} and then
\eqref{eq:ex-lattice-price-spends}. The triangle inequality yields
\[
 \|x_{j+1}\|
 \leq P_j\left(\|x_j\|
       +s_j\frac{t}{\Gamma_{j+1}}+\varepsilon_jtH_j\right).
\]
Using \(\|x_j\|\leq tI_j\), \(s_j=b_j/\kappa_{j+1}\), and the definition
of \(I_{j+1}\), the right side is
\[
 tI_{j+1}\left(
   \frac{P_jI_j}{I_{j+1}}
   +P_jb_j
   +\frac1{10}\right)
 \leq tI_{j+1}\left(\frac12+\frac16+\frac1{10}\right)
 =\frac{23}{30}tI_{j+1}\leq tI_{j+1}.
\]
Here \(P_jb_j\leq1/6\) follows from the fourth inequality in
\eqref{eq:ex-lattice-price-budget}. This proves the claim, including the
degenerate case \(t=0\).
\end{proof}

\begin{proposition}[Finite endpoint price estimate]
\label{prop:ex-lattice-trace-sharp-read}
Assume the selected history satisfies the scale inequalities
\eqref{eq:ex-lattice-price-budget} for every \(m\leq j<k\). Let a paid
comparison have coefficient vectors \(x_j\), remainder costs \(r_j\),
and marked sites at the spatial origin for \(m\leq j\leq k\). Assume
its left coefficient is zero, \(x_m=0\), and that its adjacent pairs
satisfy the local estimates \eqref{eq:ex-lattice-price-one-step} and
\eqref{eq:ex-lattice-price-spends}. Assume also the paid direction
conditions of Section~\ref{sec:reference-cone}, the room and bank
inequalities \eqref{eq:reference-bank}, and the terminal spread bound
\eqref{eq:reference-terminal} with endpoint budget \(w\geq0\) and with
the column allowances \(c_j\). Then,
for all period indices and volume indices for which the local one-step
estimates hold at every rung \(m\leq j<k\),
\begin{equation}
 \kappa_k\|x_k\|\leq4(1+c_k)w.
 \label{eq:ex-lattice-price-endpoint}
\end{equation}
If, as in the selected construction, the one-step estimates hold for all
period indices \(n\geq n_j\), uniformly in the volume index \(N\geq k\) and
in the comparison data, then the endpoint estimate holds for all
\(n\geq n_0=\max(\{0\}\cup\{n_j:m\leq j<k\})\) (the set of rungs is empty
when \(m=k\), where the estimate is \(x_k=x_m=0\)), uniformly in \(N\geq k\), in the paid
comparison and in \(w\geq0\).
\end{proposition}

\begin{proof}
We first verify that the comparison satisfies the discounted paid-cone
bound used in the recurrence. In the canonical cone estimate, choose at
rung \(j\) the worth coefficient
\(K_j=K(J(j+1))\), the column allowance equal to the selected column
allowance \(c_j\) of the history, and the bank comparison coefficient
\(M_j=K(J(j+1)+1)\). The scale-slow condition gives
\(K_j\leq K(J(j)+1)\), so the room inequality yields
\(64K_j^2\leq c_j\). The bank condition is exactly
\(100000M_j^6(1+c_{j+1})^2\leq p_j\), and monotonicity gives
\(M_j\geq K_j\geq1\). The paid direction in that cone is the adjacent
coefficient comparison with diagonal weights \(L_j^{d-r}\) and loss
\(p_j^{-3}\); this is the direction hypothesis assumed here. The zero
left coefficient and terminal spread are precisely its boundary
conditions. Therefore equation~\eqref{eq:reference-gain} gives
\begin{equation}
 \Gamma_j(r_j+\|x_j\|)\leq4(1+c_k)w
       \qquad(m\leq j\leq k).
 \label{eq:ex-lattice-price-cone}
\end{equation}

Put \(t=4(1+c_k)w\). Since \(x_m=0\), we have
\(\|x_m\|=0\leq tI_m\). Suppose inductively that
\(\|x_j\|\leq tI_j\). The paid-cone bound gives
\(r_j+\|x_j\|\leq t/\Gamma_j\) and
\(r_{j+1}+\|x_{j+1}\|\leq t/\Gamma_{j+1}\). In particular the local
remainder allowances satisfy the two bounds in Lemma~\ref{lem:ex-lattice-price-one-step}.
That lemma gives \(\|x_{j+1}\|\leq tI_{j+1}\). Finite induction reaches
\(j=k\), where \(\Gamma_k=1\), so \(I_k=1/\kappa_k\). Multiplying by \(\kappa_k\)
gives \eqref{eq:ex-lattice-price-endpoint}. The finite intersection of the
period thresholds for \(j=m,\ldots,k-1\) is still an eventual set. Each
adjacent estimate is uniform over the volume index and the specified
comparison data, so one
threshold works for all of them.
\end{proof}

Thus the endpoint estimate combines the already proved paid-cone argument
of Section~\ref{sec:reference-cone} with the one-step quantitative price
recurrence proved here. The only new induction is the finite propagation
from the zero left coefficient to the terminal rung.
